%beamer

% Define a global usable date. Must come before StyleTut
% \newcommand{\mydate}{11.11.2016}

% Comment/uncomment this line to toggle handout mode
% \newcommand{\handout}{}

% % \bigtimes abgeschrieben von http://tex.stackexchange.com/questions/14386/importing-a-single-symbol-from-a-different-font
% \DeclareFontFamily{U}{mathx}{\hyphenchar\font45}
% \DeclareFontShape{U}{mathx}{m}{n}{
%       <5> <6> <7> <8> <9> <10> gen * mathx
%       <10.95> mathx10 <12> <14.4> <17.28> <20.74> <24.88> mathx12
%       }{}
% \DeclareSymbolFont{mathx}{U}{mathx}{m}{n}
% \DeclareMathSymbol{\bigtimes}{\mathop}{mathx}{161}

\RequirePackage{xcolor}

\def\9{\square}
%\def\9{\blank}

% f"ur Aussagenlogik
\colorlet{alcolor}{blue}
\RequirePackage{tikz}
\usetikzlibrary{arrows.meta}
\newcommand{\alimpl}{\mathrel{\tikz[x={(0.1ex,0ex)},y={(0ex,0.1ex)},>={Classical TikZ Rightarrow[]}]{\draw[alcolor,->,line width=0.7pt,line cap=round] (0,0) -- (15,0);\path (0,-6);}}}
\newcommand{\aleqv}{\mathrel{\tikz[x={(0.1ex,0ex)},y={(0ex,0.1ex)},>={Classical TikZ Rightarrow[]}]{\draw[alcolor,<->,line width=0.7pt,line cap=round] (0,0) -- (18,0);\path (0,-6);}}}
\newcommand{\aland}{\mathbin{\raisebox{-0.6pt}{\rotatebox{90}{\texttt{\color{alcolor}\char62}}}}}
\newcommand{\alor}{\mathbin{\raisebox{-0.8pt}{\rotatebox{90}{\texttt{\color{alcolor}\char60}}}}}
%\newcommand{\ali}[1]{_{\mathtt{\color{alcolor}#1}}}
\newcommand{\alv}[1]{\mathtt{\color{alcolor}#1}}
\newcommand{\alnot}{\mathop{\tikz[x={(0.1ex,0ex)},y={(0ex,0.1ex)}]{\draw[alcolor,line width=0.7pt,line cap=round,line join=round] (0,0) -- (10,0) -- (10,-4);\path (0,-8) ;}}}
\newcommand{\alP}{\alv{P}} %ali{#1}}
%\newcommand{\alka}{\negthinspace\hbox{\texttt{\color{alcolor}(}}}
\newcommand{\alka}{\negthinspace\text{\texttt{\color{alcolor}(}}}
%\newcommand{\alkz}{\texttt{\color{alcolor})}}\negthinspace}
\newcommand{\alkz}{\text{\texttt{\color{alcolor})}}\negthinspace}
\newcommand{\AAL}{A_{AL}}
\newcommand{\LAL}{\hbox{\textit{For}}_{AL}}
\newcommand{\AxAL}{\hbox{\textit{Ax}}_{AL}}
\newcommand{\AxEq}{\hbox{\textit{Ax}}_{Eq}}
\newcommand{\AxPL}{\hbox{\textit{Ax}}_{PL}}
\newcommand{\AALV}{\hbox{\textit{Var}}_{AL}}
\newcommand{\MP}{\hbox{\textit{MP}}}
\newcommand{\GEN}{\hbox{\textit{GEN}}}
\newcommand{\W}{\ensuremath{\hbox{\textbf{w}}}\xspace}
\newcommand{\F}{\ensuremath{\hbox{\textbf{f}}}\xspace}
\newcommand{\WF}{\ensuremath{\{\W,\F\}}\xspace}
\newcommand{\val}{\hbox{\textit{val}}}
\newcommand{\valDIb}{\val_{D,I,\beta}}

\newcommand*{\from}{\colon}

% die nachfolgenden Sachen angepasst an cmtt
\newlength{\ttquantwd}
\setlength{\ttquantwd}{1ex}
\newlength{\ttquantht}
\setlength{\ttquantht}{6.75pt}
\def\plall{%
  \tikz[line width=0.67pt,line cap=round,line join=round,baseline=(B),alcolor] {
    \draw (-0.5\ttquantwd,\ttquantht) -- node[coordinate,pos=0.4] (lll){} (-0.25pt,-0.0pt) -- (0.25pt,-0.0pt) -- node[coordinate,pos=0.6] (rrr){} (0.5\ttquantwd,\ttquantht);
    \draw (lll) -- (rrr);
    \coordinate (B) at (0,-0.35pt);
  }%
}
\def\plexist{%
  \tikz[line width=0.67pt,line cap=round,line join=round,baseline=(B),alcolor] {
    \draw (-0.9\ttquantwd,\ttquantht) -- (0,\ttquantht) -- node[coordinate,pos=0.5] (mmm){} (0,0) --  (-0.9\ttquantwd,0);
    \draw (mmm) -- ++(-0.75\ttquantwd,0);
    \coordinate (B) at (0,-0.35pt);
  }\ensuremath{\,}%
}
\let\plexists=\plexist
\newcommand{\NT}[1]{\ensuremath{\langle\mathrm{#1} \rangle}}

\newcommand{\CPL}{\text{\itshape Const}_{PL}}
\newcommand{\FPL}{\text{\itshape Fun}_{PL}}
\newcommand{\RPL}{\text{\itshape Rel}_{PL}}
\newcommand{\VPL}{\text{\itshape Var}_{PL}}
\newcommand{\ATer}{A_{\text{\itshape Ter}}}
\newcommand{\ARel}{A_{\text{\itshape Rel}}}
\newcommand{\AFor}{A_{\text{\itshape For}}}
\newcommand{\LTer}{L_{\text{\itshape Ter}}}
\newcommand{\LRel}{L_{\text{\itshape Rel}}}
\newcommand{\LFor}{L_{\text{\itshape For}}}
\newcommand{\NTer}{N_{\text{\itshape Ter}}}
\newcommand{\NRel}{N_{\text{\itshape Rel}}}
\newcommand{\NFor}{N_{\text{\itshape For}}}
\newcommand{\PTer}{P_{\text{\itshape Ter}}}
\newcommand{\PRel}{P_{\text{\itshape Rel}}}
\newcommand{\PFor}{P_{\text{\itshape For}}}

\newcommand{\plka}{\alka}
\newcommand{\plkz}{\alkz}
%\newcommand{\plka}{\plfoo{(}}
%\newcommand{\plkz}{\plfoo{)}}
\newcommand{\plcomma}{\hbox{\texttt{\color{alcolor},}}}
\newcommand{\pleq}{{\color{alcolor}\,\dot=\,}}

% MODIFIED (DJ)
% previously: \newcommand{\plfoo}[1]{\mathtt{\color{alcolor}#1}}
\newcommand{\plfoo}[1]{\texttt{\color{alcolor}#1}}

\newcommand{\plc}{\plfoo{c}}
\newcommand{\pld}{\plfoo{d}}
\newcommand{\plf}{\plfoo{f}}
\newcommand{\plg}{\plfoo{g}}
\newcommand{\plh}{\plfoo{h}}
\newcommand{\plx}{\plfoo{x}}
\newcommand{\ply}{\plfoo{y}}
\newcommand{\plz}{\plfoo{z}}
\newcommand{\plR}{\plfoo{R}}
\newcommand{\plS}{\plfoo{S}}

\newcommand{\bv}{\mathrm{bv}}
\newcommand{\fv}{\mathrm{fv}}

%\newcommand{\AxAL}{\hbox{\textit{Ax}}_{AL}}
%\newcommand{\AALV}{\hbox{\textit{Var}}_{AL}}

%\renewcommand{\#}[1]{\literal{#1}}
\newcommand{\A}{\mathcal{A}}
\newcommand{\Adr}{\text{Adr}}
\newcommand{\ar}{\mathrm{ar}}
\newcommand{\ascii}[1]{\literal{\char#1}}
%\newcommand{\assert}[1]{\text{/\!\!/\ } #1}
\newcommand{\assert}[1]{\colorbox{black!7!white}{\ensuremath{\{\;#1\;\}}}}
\newcommand{\Assert}[1]{$\langle$\textit{#1}$\rangle$}
\newcommand{\B}{\mathcal{B}}
\newcommand{\bfmod}{\mathbin{\kw{ mod }}}
\newcommand{\bb}{{\text{bb}}}
\def\bottom{\hbox{\small$\pmb{\bot}$}}
\newcommand{\card}[1]{|#1|}
%\newcommand{\cod}{\mathop{\text{cod}}}  % ist in thwmathabbrevs
\newcommand{\Conf}{\mathcal{C}}
\newcommand{\define}[1]{\emph{#1}}
%\renewcommand{\dh}{d.\,h.\@\xspace}
%\newcommand{\Dh}{D.\,h.\@\xspace}
%\newcommand{\engl}[1]{engl.\xspace\emph{#1}}
\newcommand{\eps}{\varepsilon}
%\newcommand{\evtl}{evtl.\@\xspace}
\newcommand{\fbin}{\text{bin}}
\newcommand{\finv}{\text{inv}}
\newcommand{\fnum}{\text{num}}
\newcommand{\fNum}{{\text{Num}}}
\newcommand{\frepr}{\text{repr}}
\newcommand{\fRepr}{\text{Repr}}
\newcommand{\fZkpl}{\text{Zkpl}}
\newcommand{\fLen}{\text{Len}}
\newcommand{\fsem}{\text{sem}}
\providecommand{\fspace}{\mathord{\text{space}}}
\providecommand{\fSpace}{\mathord{\text{Space}}}
\providecommand{\ftime}{\mathord{\text{time}}}
\providecommand{\fTime}{\mathord{\text{Time}}}
\newcommand{\fTrans}{\text{Trans}}
\newcommand{\fVal}{\text{Val}}

% MODIFIED (DJ)
\newcommand{\Val}{\text{Val}}

%\def\G{\mathbb{Z}}
\newcommand{\HT}[1]{\normalfont\textsc{HT-#1}}
\newcommand{\htr}[3]{\{#1\}\;#2\; \{#3\}}
\newcommand{\Id}{\text{I}}
%\newcommand{\ie}{i.\,e.\@\xspace}
\newcommand{\instr}[2]{\texttt{#1}\ \textit{#2}}
\newcommand{\Instr}[2]{\texttt{#1}\ \textrm{#2}}
\newcommand{\instrr}[3]{\texttt{#1}\ \textit{#2}\texttt{(#3)}}
\newcommand{\Instrr}[3]{\texttt{#1}\ \textrm{#2}\texttt{(#3)}}
\newcommand{\io}{\!\mid\!}
\usepackage{KITcolors}
\newcommand{\literal}[1]{\hbox{\textcolor{blue!95!white}{\textup{\texttt{\scalebox{1.11}{#1}}}}}}
%\newcommand{\literal}[1]{\hbox{\textcolor{KITblue!80!black}{\textup{\texttt{#1}}}}}
\def\kasten#1{\leavevmode\literal{\setlength{\fboxsep}{1pt}\fbox{\vrule  width 0pt height 1.5ex depth 0.5ex #1}}}
\newcommand{\kw}[1]{\ensuremath{\mathbf{#1}}}
\newcommand{\lang}[1]{\ensuremath{\langle#1\rangle}}
%\newcommand{\maw}{m.\,a.\,w.\@\xspace}
%\newcommand{\MaW}{M.\,a.\,w.\@\xspace}
\newcommand{\mdefine}[2][FOOBAR]{\define{#2}\def\foobar{FOOBAR}\def\optarg{#1}\ifx\foobar\optarg\def\optarg{#2}\fi\graffito{\optarg}}
\newcommand{\meins}{\rotatebox[origin=c]{180}{1}}
\newcommand{\Mem}{\text{Mem}}
\newcommand{\memread}{\text{memread}}
\newcommand{\memwrite}{\text{memwrite}}
\providecommand{\meta}[1]{\ensuremath{\langle}\textit{#1}\ensuremath{\rangle}}
%\newcommand{\N}{\mathbb{N}}
\newcommand{\NP}{\mathbf{NP}}
\newcommand{\Nadd}{N_{\text{add}}}
\newcommand{\Nmult}{N_{\text{mult}}}
% MODIFIED (DJ): added \!, mathcal{O}
\newcommand{\Oh}[1]{\mathcal{O}\!\left(#1\right)}
\newcommand{\Om}[1]{\Omega\!\left(#1\right)}
\newcommand{\personname}[1]{\textsc{#1}}
\newcommand{\regname}[1]{\texttt{#1}}
\newcommand{\mima}{\textsc{Mima}\xspace}
\newcommand{\mimax}{\textsc{Mima-X}\xspace}

\def\Pclass{\text{\bfseries P}}
\def\PSPACE{\text{\bfseries PSPACE}}

\newcommand{\SPush}{\text{push}}
\newcommand{\SPop}{\text{pop}}
\newcommand{\SPeek}{\text{peek}}
\newcommand{\STop}{\text{top}}
\newcommand{\STos}{\text{\itshape tos}}
\newcommand{\SBos}{\text{\itshape bos}}

%\newcommand{\R}{\mathbb{R}}
\newcommand{\Rnullplus}{\R_0^{+}}
\newcommand{\Rplus}{\R_{+}}
\newcommand{\resp}{resp.\@\xspace}
\newcommand{\Sem}{\text{Sem}}
\newcommand{\sgn}{\mathop{\text{sgn}}}
\newcommand{\sqbox}{\mathop{\raisebox{-6.2pt}{\hbox{\hbox to 0pt{$^{^{\sqcap}}$\hss}$^{^{\sqcup}}$}}}}
\newcommand{\sqleq}{\sqsubseteq}
\newcommand{\sqgeq}{\sqsupseteq}
% MODIFIED (DJ): added \!
\newcommand{\Th}[1]{\Theta\!\left(#1\right)}
%\newcommand{\usw}{usw.\@\xspace}
\newcommand{\V}[1]{\hbox{\textit{#1}}}
\newcommand{\x}{\times}
\newcommand{\ZK}{\mathbb{K}}
%\newcommand{\Z}{\mathbb{Z}}
\newcommand{\zB}{z.\,B.\@\xspace}
\newcommand{\ZB}{Z.\,B.\@\xspace}
% \newcommand{\bb}{{\text{bb}}}
% \def\##1{\hbox{\textcolor{darkblue}{\texttt{#1}}}}
% \def\A{\mathcal{A}}
% \newcommand{\0}{\#0}
% \newcommand{\1}{\#1}
% \newcommand{\Obj}{\text{Obj}}
% \newcommand{\start}{\mathop{\text{start}}}
% \newcommand{\compactlist}{\addtolength{\itemsep}{-\parskip}}
% \newcommand{\fval}{\text{val}}
% \newcommand{\lang}[1]{\ensuremath{\langle#1\rangle}}
% \newcommand{\io}{\!\mid\!}
% \def\sqbox{\mathop{\raisebox{-6.2pt}{\hbox{\hbox to 0pt{$^{^{\sqcap}}$\hss}$^{^{\sqcup}}$}}}}
% \def\sqleq{\sqsubseteq}
% \def\sqgeq{\sqsupseteq}
\def\Td{T_{\overline{d}}}
% \newcommand{\csym}[1]{\ensuremath{\#{c}_{\#{\hbox{\scriptsize #1}}}}}
% \newcommand{\F}{\ensuremath{\mathcal{F}}}
% \newcommand{\fsym}[2]{\ensuremath{\#{f}^{\#{\hbox{\scriptsize #1}}}_{\#{\hbox{\scriptsize #2}}}}}
% \newcommand{\rsym}[2]{\ensuremath{\#{R}^{\#{\hbox{\scriptsize #1}}}_{\#{\hbox{\scriptsize #2}}}}}
% \newcommand{\xsym}[1]{\ensuremath{\#{x}_{\#{\hbox{\scriptsize #1}}}}}
% \newcommand{\I}{\mathcal{I}}
% ********************************************************************

\usepackage{../TutTexbib/thwregex}

\newcommand{\nM}{\mathbb{M}}
\newcommand{\nR}{\mathbb{R}}
\newcommand{\nN}{\mathbb{N}}
\newcommand{\nZ}{\mathbb{Z}}
\newcommand{\nQ}{\mathbb{Q}}
\newcommand{\nB}{\mathbb{B}}
\newcommand{\nC}{\mathbb{C}}
\newcommand{\nK}{\mathbb{K}}
\newcommand{\nF}{\mathbb{F}}
\newcommand{\nG}{\mathbb{G}}
\newcommand{\nullel}{\mathcal{O}}
\newcommand{\einsel}{\mathds{1}}
\newcommand{\nP}{\mathbb{P}}
\newcommand{\Pot}{\mathcal{P}}
\renewcommand{\O}{\text{O}}

%\newcommand{\literal}[1]{\textcolor{blue}{\texttt{#1}}}
\newcommand{\realTilde}{\textasciitilde \ }
\newcommand{\setsize}[1]{\; \mid #1 \mid \; }

\definecolor{myRed}{RGB}{255,75,20}
\colorlet{myGreen}{KITpalegreen}

\newcounter{tfqtempcount}
\newcommand{\truefalseQ}[4]{
	\setcounter{tfqtempcount}{#1}
	\addtocounter{tfqtempcount}{1}
	\truefalseQuestion{#1}{\value{tfqtempcount}}{#2}{#3}{#4}
}
\newcommand{\truefalseQuestion}[5]{\item<#1-|handout:#1-> \color<#2-|handout:#2->{#3} #4 \qquad \visible<#2-|handout:#2->{#5}}

\newcommand{\boder}{\ensuremath{\textcolor{blue}{\vee}}\text{} }
\newcommand{\bund}{\ensuremath{\textcolor{blue}{\wedge}}\text{}}
\newcommand{\bimp}{\ensuremath{\textcolor{blue}{\to}}\text{}}
\newcommand{\bnot}{\ensuremath{\textcolor{blue}{\neg}}\text{}}
\newcommand{\bone}{\ensuremath{\textcolor{blue}{1}}\text{}}
\newcommand{\bzero}{\ensuremath{\textcolor{blue}{0}}\text{}}
\newcommand{\bleftBr}{\ensuremath{\textcolor{blue}{(}}\text{}}
\newcommand{\brightBr}{\ensuremath{\textcolor{blue}{)}}\text{}}

\newcommand{\summe}[2]{\sum\limits_{#1}^{#2}}
\newcommand{\limes}[1]{\lim\limits_{#1}}

\newcommand{\numpp}{\advance \value{weeknum} by -2 \theweeknum \advance \value{weeknum} by 2}
\newcommand{\nump}{\advance \value{weeknum} by -1 \theweeknum \advance \value{weeknum} by 1}

\newcommand{\mycomment}[1]{}

\newcounter{abc}
\newenvironment{alist}{
  \begin{list}{(\alph{abc})}{
      \usecounter{abc}\setlength{\leftmargin}{8mm}\setlength{\labelsep}{2mm}
    }
}{\end{list}}

\renewcommand{\mod}{\text{ \bf{mod} }}
\renewcommand{\div}{\text{ \bf{div} }}

%\newcommand{\sgn}{\text{sgn}}

% Toggle Handout mode by including this line before including style_tut:
% \newcommand{\handout}{}

%% Beamer-Klasse im korrekten Modus
\ifdefined \handout
\documentclass[handout]{beamer} % Handout mode
\else
\documentclass{beamer}
\fi


% Das ist der KIT-Stil
%\usepackage{../TutTexbib/beamerthemekit}
\usepackage[deutsch,titlepage0]{../TutTexbib/KIT/beamerthemeKITmod}
\TitleImage[width=\titleimagewd]{../figures/titlepage.jpg}
%\usetheme[deutsch,titlepage0]{KIT}

% Include PDFs
\usepackage{pdfpages}

% Libertine font (Original GBI font)
\usepackage{libertine}
%\renewcommand*\familydefault{\sfdefault}  %% Only if the base font of the document is to be sans serif

% Nicer math symbols
\usepackage{eulervm}
%\usepackage{mathpazo}

%% Deutsche Silbentrennung und Beschriftungen
\usepackage[ngerman]{babel}

%% UTF-8-Encoding
\usepackage[utf8]{inputenc}

\usepackage{csquotes}

%% Bibliotheken für viele mathematische Symbole
\usepackage{amsmath, amsfonts, amssymb}

%% Anzeigetiefe für Inhaltsverzeichnis: 1 Stufe
\setcounter{tocdepth}{1}

%% Schönere Schriften
\usepackage[TS1,T1]{fontenc}

%% Tabellen
\usepackage{array}
\usepackage{multicol}

%% Bibliothek für Graphiken
\usepackage{graphicx}

%% der wird sowieso in jeder Datei gesetzt
%%\graphicspath{{../figures/}}

%% Hyperlinks
\usepackage{hyperref}
% I don't know why, but this works and only includes sections and NOT subsections in the pdf-bookmarks.
\hypersetup{bookmarksdepth=subsection} 

%\usepackage{lmodern}
\usepackage{colortbl}
\usepackage[absolute,overlay]{textpos}
\usepackage{listings}
\usepackage{forloop}
%\usepackage{algorithmic} % PseudoCode package 

\usepackage{tikz}
\usetikzlibrary{matrix}
\usetikzlibrary{arrows.meta}
\usetikzlibrary{automata}
\usetikzlibrary{tikzmark}

% Needed for gbi-macros
\usepackage{xspace}

%%%%%%%%%%%% INHALT %%%%%%%%%%%%%%%%

%% Wochennummer
\newcounter{weeknum}

%% Titelinformationen
\title[GBI Tutorium, Woche \theweeknum]{Grundbegriffe der Informatik \\ Tutorium 36}
\subtitle{Termin \theweeknum \ | \mydate \\ Thassilo Helmold}
\author[Thassilo Helmold]{Thassilo Helmold}
\institute{KIT -- Karlsruher Institut für Technologie}
\date{\mydate}

%% Titel einfügen
\newcommand{\titleframe}{\frame{\titlepage}}

%% Alles starten mit \starttut{X}
\newcommand{\starttut}[1]{\setcounter{weeknum}{#1}\titleframe\frame{\frametitle{Inhalt}\tableofcontents} \AtBeginSection[]{%
\begin{frame}
	\tableofcontents[currentsection]
\end{frame}\addtocounter{framenumber}{-1}}}


\newcommand{\framePrevEpisode}{
	\begin{frame}
		\centering
		\textbf{In the previous episode of GBI...}
	\end{frame}
}

%% Legacy: Lastframe. Not for further usage!
\newcommand{\lastframe}[4]{\xkcdframe{#1}{#2}{#3}{#4}{}}
\newcommand{\xkcdframe}[5]{
	\frame[plain]{
		\vspace{-#2pt}
		\begin{figure}[H]
			\centering
			\includegraphics[scale=#1]{#3}
			\vspace{-7pt}
			\caption{ \texttt{\url{#4}} }
			\vspace{5pt}
			#5
		\end{figure} 
	}
}

%% Wörter
\newcommand{\code}[1]{$\mathbf{#1}$}

%% Sterne

\newcounter{starsc}
\newcommand{\stars}[1]{
	\hfill
	\begin{minipage}{100px}
		\forloop{starsc}{0}{\value{starsc} < #1}%
		{%
			\includegraphics[scale=0.05]{star-full.pdf} \hspace*{1px}
		}%
		\forloop{starsc}{\value{starsc}}{\value{starsc} < 5}%
		{%
			\includegraphics[scale=0.05]{star-empty.pdf} \hspace*{1px}
		}
		\vspace*{2px}
	\end{minipage}
}

% No stars for me...
\renewcommand{\stars}[1]{}

\newcommand{\slideThanks}{
	\begin{frame}
		\frametitle{Danksagung}
		\begin{block}{}
			Dieser Foliensatz basiert in Teilen auf Folien von:\\[1em]
			Philipp Basler \\
			Nils Braun \\
			Dominik Doerner \\
			Ou Yue \\
		\end{block}
	\end{frame}
}

%% Verbatim
\usepackage{moreverb}

\graphicspath{{../figures/}}

\begin{document}
\starttut{3}

\begin{frame}
	\frametitle{Übungsblätter}
	\pause
	... gibt es natürlich noch nicht ;-) \\[1em]
	\pause
	Noch ein paar spezielle Hinweise:
	\begin{itemize}
		\item Text in Schwarz oder Blau (oder andere dunkle Farbe $\notin \{\text{Rot}, \text{Grün}\}$ )
		\item Folgende Farben sind für Kennzeichnungen zugelassen:\\
		Blau, Schwarz, Dunkelgrün, Orange, Gelb.\\
		Andere Farben, insbesondere Rot, sind nicht zugelassen \\
		(und es gilt Blau = Lila, Gelb = Hellgrün, ...)
		\item Links und rechts bitte jeweils einen Rand von mind. 1,5 cm frei lassen.
		\item Aufgabennummer deutlich kenntlich machen
		\item Für (sehr) unsaubere Abgaben ist Punktabzug möglich!
	\end{itemize}
	
	Und ganz wichtig: Frühzeitig anfangen! ÜB 2 sowie folgende ÜB sind \textbf{doppelt so umfangreich} wie ÜB 1 (und geben entsprechend viele Punkte)!
\end{frame}

\framePrevEpisode

\subsection{Wahr oder Falsch?}
\begin{frame}
	\begin{block}{Wahr oder Falsch?}
		\begin{itemize}
			\FalseQuestionE{Eine Funktion muss linkseindeutig sein}{Richtig: Eine Funktion muss rechtseindeutig und linkstotal sein}
			\TrueQuestionE{Eine injektive Funktion ist linkseindeutig}{}
			\TrueQuestionE{Eine surjektive Funktion ist rechtstotal}{}
			\FalseQuestion{Jede Relation ist eine Funktion}
			\TrueQuestionE{Jede Funktion ist eine Relation}{}
		\end{itemize}
	\end{block}
\end{frame}

\begin{frame}
	\begin{block}{Wahr oder Falsch?}
	\begin{itemize}
		\FalseQuestionE{$aaba \in \{a,b\}^2\times\{a,b\}^2$}{Aber: $(\text{aa}, \text{ba}) \in \{a,b\}^2\times\{a,b\}^2$}
		\FalseQuestionE{Ich benutze für jeden Beweis Induktion.}{Mit einem Schraubenschlüssel bekommt man keinen Nagel in die Wand.}
		\FalseQuestionE{$\setsize{\{ \varepsilon \}} = 0$}{\{ \varepsilon \} \neq \emptyset$}
	\end{itemize}
	\end{block}
\end{frame}

\frame{
	\frametitle{Zum Aufwärmen: Logikrätsel}
	\vspace{-2pt}
	\begin{figure}[H]
		\centering
		\only<1|handout:1>{\includegraphics[scale=1]{xkcd/logic_boat_problem}}
		\only<2|handout:2>{\includegraphics[scale=0.35]{xkcd/logic_boat_sol}}
		\vspace{-7pt}
		\caption{ \texttt{\url{http://xkcd.com/1134/}} }
	\end{figure} 
}

\section{Aussagenlogik}

\thasse{ % No time this year...
	\frame{
		\frametitle{Ein bekanntes Logikrätsel}
		\vspace{-2pt}
		\begin{figure}[H]
			\centering
			\only<1|handout:1>{\includegraphics[scale=1]{xkcd/logic_boat_problem}}
			\only<2|handout:2>{\includegraphics[scale=0.35]{xkcd/logic_boat_sol}}
			\vspace{-2pt}
			\url{http://xkcd.com/1134/}
		\end{figure} 
	}
}

\begin{frame}{Aussagenlogik}
	Aussagen sind Sätze, die entweder \textbf{wahr} oder \textbf{falsch} sind.\\
	Deren Wahrheitswert muss dabei nicht unbedingt bekannt oder \enquote{tatsächlich ermittelbar} sein.
	
	\pause
	\begin{Beispiel}
		\begin{itemize}
			\item „$ 1 + 1 = 2 $“ ist eine Aussage. Sie ist wahr.
			\item \enquote{Es gibt nur endlich viele Primzahlen.} ist eine Aussage. Sie ist falsch.
			\pause
			\item Die Goldbachsche Vermutung (Jede gerade Zahl größer 2 ist Summe zweier Primzahlen) \pause ist eine Aussage. Ihr Wahrheitswert ist unbekannt.
			\pause
			\item „Die Welt wird am 11.11.11111 untergehen.“ \pause ist auch eine Aussage. Wir werden ihren Wahrheitswert aber wohl niemals ermitteln können.
			\pause
			\item \enquote{Gelb} ist keine Aussage.
			\pause
			\item \enquote{Dieser Satz ist falsch.} \pause ist keine Aussage. Dem Satz kann offensichtlich kein Wahrheitswert zugeordnet werden.
		\end{itemize}
	\end{Beispiel}
\end{frame}

\begin{frame}{Zwei Grundsätze der Aussagenlogik}
	
	\begin{itemize}
		\pause
		\item[1.] Jede Aussage ist \textbf{entweder} wahr \textbf{oder} falsch.\\
		Wir schreiben im Folgenden $\BB= \{\W, \F\}$
		\pause
		\item[2.] Der Wahrheitswert einer \textbf{zusammengesetzten} Aussage ist durch die
		Wahrheitswerte der \textbf{Teilaussagen} eindeutig festgelegt. \\
		\enquote{$2 + 2 = 5 \; \bimp \; \text{Pinguine können fliegen}$} ist \textbf{wahr}.\\[0.2em]
		%TODO
		%Ex falso quodlibet
	\end{itemize}

	\impl Weg vom Inhalt und betrachte stattdessen \textbf{Aussagevariablen}.
\end{frame}

\begin{frame}{Syntax}
	\begin{Definition}
		$Var_{AL}$ ist die Menge aller \textbf{Aussagevariablen}. \\
		Wir bilden aussagenlogische Formeln als Wörter über dem Alphabet $A_{AL} = \{ \bleftBr, \brightBr, \bnot, \bund, \boder, \bimp \} \cup Var_{AL}$ 
	\end{Definition}
	\pause
	\begin{Definition}
		$For_{AL}$ ist die Menge aller \textbf{syntaktisch korrekten Formeln} über $Var_{AL}$.\\
		\medskip
		(VL: induktive Definition über Konstruktionsabbildungen) \\
		Klammern, die überflüssig sind, können weg. Wir lesen es dann wie gewohnt (\enquote{$\bnot$} vor \enquote{ $\bund$} vor \enquote{ $\boder$} vor \enquote{ $\bimp$}  {\small \#PunktVorStrichUndSo}).
	\end{Definition}
	\pause
	\begin{Beispiel}
		$$Var_{AL} = \{\alA, \alB, \alC\}$$
		$$For_{AL} = \set{\alka \alA  \bimp \alB \alkz \boder \bnot \alB \bund \alC, \ \alC \bimp \alka \bnot \alC \alkz, ...}$$ 
	\end{Beispiel}
\end{frame}

\begin{frame}{Semantik}
	Die Semantik einer aussagenlogischen Formel wird durch ihre \textbf{Auswertung} bestimmt.\\
	Hierbei werden den Operatorsymbolen aus $A_{AL}$ boolesche Funktionen zugeordnet.
	
	\pause
	\begin{Definition}
		Eine \textbf{boolesche Funktion} ist eine Abbildung der Form
		$f \from \BB^n \functionto \BB$.
	\end{Definition}

	\pause
	\begin{Beispiel}
		\enquote{Handelsübliche} boolesche Funktionen sind  $b_{\bnot}$,
		$b_{\bund}$, $b_{\boder}$ und $b_{\bimp}$. 
	\end{Beispiel}
\end{frame}

\begin{frame}{Semantik}
	\let\w\W
	\let\f\F
	\begin{center}
		\begin{huge}
			$$\only<1-2|handout:1>{\bnot}\only<3-4|handout:2>{\bund}\only<5-6|handout:3>{\boder}\only<7-8|handout:4>{\bimp}$$
		\end{huge}
		Wahrheitstabelle: \\
		\medskip
		\only<1-2|handout:1>{
			\begin{tabular}{|c|c|}
				\hline 
				$A$ &  $\bnot A$ \\
				\hline
				\w & \uncover<2|handout:1>{\f} \\
				\hline
				\f & \uncover<2|handout:1>{\w} \\
				\hline
			\end{tabular}
		}
		\only<3-4|handout:2>{ 
			\begin{tabular}{|c|c|c|}
				\hline 
				$A$ & $B$ & $A \mathop{\bund} B$ \\
				\hline
				\w & \w & \uncover<4|handout:2>{\w} \\
				\hline
				\w & \f & \uncover<4|handout:2>{\f} \\
				\hline
				\f & \w & \uncover<4|handout:2>{\f} \\
				\hline
				\f & \f & \uncover<4|handout:2>{\f} \\
				\hline
			\end{tabular}
		}
		\only<5-6|handout:3>{ 
			\begin{tabular}{|c|c|c|}
				\hline 
				$A$ & $B$ & $A \mathop{\boder} B$ \\
				\hline
				\w & \w & \uncover<6|handout:3>{\w} \\
				\hline
				\w & \f & \uncover<6|handout:3>{\w} \\
				\hline
				\f & \w & \uncover<6|handout:3>{\w} \\
				\hline
				\f & \f & \uncover<6|handout:3>{\f} \\
				\hline
			\end{tabular}
		}
		\only<7-8|handout:4>{
			\begin{tabular}{|c|c|c|}
				\hline 
				$A$ & $B$ & $A \mathop{\bimp} B$ \\
				\hline
				\w & \w & \uncover<8|handout:4>{\w} \\
				\hline
				\w & \f & \uncover<8|handout:4>{\f} \\
				\hline
				\f & \w & \uncover<8|handout:4>{\w} \\
				\hline
				\f & \f & \uncover<8|handout:4>{\w} \\
				\hline
			\end{tabular}
		}
	\end{center}
\end{frame}

\begin{frame}{Semantik}
	Wenn wir den Wahrheitswert einer \textbf{zusammengesetzten} Aussage bestimmen wollen, brauchen wir die Werte der \textbf{verwendeten Aussagenvariablen}. \\
	\begin{Definition}
		Sei $V \subseteq Var_{AL}$ die Menge der verwendeten Aussagevariablen.\\
		Eine Funktion $I \from V \functionto \BB$ bezeichnet man als \textbf{Interpretation}.
	\end{Definition}
	
	

\end{frame}

\begin{frame}{Auswertung}
	\begin{Definition}
		Sei $I$ eine Interpretation. Dann liefert die Abbildung $val_I(A)$ den Wahrheitswert zu einer aussagenlogischen Formel $A$. \\
		\medskip
		Diese Funktion nennen wir \textbf{Auswertung} und berechnen sie schrittweise, z.~B.:
	\end{Definition}
	
	\begin{align*}
	&val_I(\alka\alA \bimp \alB\alkz \boder \bnot \alB)  \\
	\visible<2->{= \;&b_{\boder} (val_I(\alA\bimp \alB), val_I(\bnot \alB)) \\}
	\visible<3->{= \;&b_{\boder} (b_{\bimp} (val_I(\alA), val_I(\alB)), b_{\bnot}(val_I(\alB))) \\}
	\visible<4->{= \;&b_{\boder} (b_{\bimp} (I(\alA), I(\alB)), b_{\bnot}(I(\alB))) \\}
	\end{align*}
\end{frame}

\begin{frame}{Auswertung}
	\begin{Definition}
		Eine \textbf{Tautologie} ist eine aussagenlogische Formel, bei der für alle möglichen Interpretationen $I$ gilt: $val_I(A) = \W$. \\
		\impl Schreibweise: \quad $\models A$ \\
		\medskip
		
		\pause
		Zwei aussagenlogische Formeln $A$ und $B$ heißen \textbf{äquivalent}, wenn für jede beliebige Interpretation $I$
		$$val_I(A) = val_I(B)$$
		gilt. \\
		\impl Schreibweise: \quad $A \equiv B$
	\end{Definition}
\end{frame}

% TODO
\begin{frame}[t]{Auswertung}
	\only<all:1>{\includegraphics[scale=0.65]{al_uebung_1}\\[2em]
	Quelle: GBI-Übung 2015/2016}
	\only<all:2>{\includegraphics[scale=0.65]{al_uebung_2}}
	\only<all:3>{\includegraphics[scale=0.65]{al_uebung_3}}
	\only<all:4>{\includegraphics[scale=0.65]{al_uebung_4}}
	\only<all:5>{\includegraphics[scale=0.65]{al_uebung_5}}
\end{frame}

\begin{frame}{Semantik}
	
	Möchte man eine AL-Formel für alle möglichen Interpretationen auswerten, so macht man dies meist in Form einer \textbf{Wahrheitstabelle}.
	
	\begin{block}{Aufgabe}
		Gegeben seien die Formeln
		$$ F_1 = \left(\left(\left(\alB \bimp \alA \right) \boder \alB \right) \bimp (\bnot \alA)\right) \bund \alB$$
		und
		$$F_2 = \bnot \alA \bund \alB$$
		Stellt die Wahrheitstabellen von $F_1$ und $F_2$ auf. Sind die beiden Formeln äquivalent?
	\end{block}
\end{frame}

\begin{frame}{Lösung}
	Für die Formel $F_1$:
	\begin{table}[H]
	\centering
	\begin{tabular}{|*{6}{c|}}
	\hline
	$\alA$ & $\alB$ & $\alB \bimp \alA$ &  $\dots \boder \alB$ & $\dots \bimp \bnot \alA$ & $\dots \bund \alB$ \pause \\ \hline
	\W & \W & \W & \W & \F & \F  \\ \hline \pause
	\W &  \F & \W & \W & \F & \F  \\ \hline \pause
	\F & \W & \F & \W & \W & \W  \\ \hline \pause
	\F & \F & \W & \W & \W & \F \\ \hline
	\end{tabular}
	\end{table}
\end{frame}

\begin{frame}{Lösung}
	Für die Formel $F_2$:
	\begin{table}[h!]
	\centering
	\begin{tabular}{|*{3}{c|}}
	\hline
	$\alA$ & $\alB$ & $\bnot \alA \bund \alB$  \\ \hline \pause
	\W & \W & \F \\  \hline \pause
	\W & \F & \F \\  \hline \pause
	\F & \W & \W \\  \hline \pause
	\F & \F & \F \\ \hline
	\end{tabular}
	\end{table}
	Also sind die beiden Formeln äquivalent: $$F_1 \equiv F_2$$
\end{frame}


%\input{../Bloecke/VollstaendigeInduktion.tex}

\section{Formale Sprachen}

\begin{frame}{Rückblick}
	Sei $A = \{\word A, \word B, ..., \word Z, \word a, \word b, ..., \word z, \word\sp\}$ ein Alphabet.\\
	\pause
	Dann enthält $A^*$ alle Wörter, die man mit Zeichen aus $A$ bilden kann. Aber nicht jedes dieser Wörter ist auch sinnvoll.\\[1em]
	
	\enquote{\word{\thassedaniel{egnarts\sp si\sp efiL}{retsinnaL\sp nosrO}}} ist kein sinnvolles Wort. \pause Oder? \\[1em]
	\pause
	\impl Hängt vom Kontext ab!\\
	
\end{frame}

\begin{frame}{Formale Sprachen}
		\begin{Definition}
			Sei $A$ ein Alphabet.\\
			Eine \textbf{formale Sprache} $L$ ist eine Teilmenge von $A^*$.
		\end{Definition}
		\pause
		\vspace{10pt}
		Durch eine formale Sprache geben wir an, welche der möglichen Wörter über $A$ wir in einem bestimmten Kontext als „wünschenswert“ ansehen.\\
		\pause
		Formale Sprachen werden häufig nicht direkt, sondern über Bildungsvorschriften angegeben.
		
		\pause
		\begin{Beispiel}
			$A = \{\word 0, \word 1\}$ \\
			$L = \{ w \in A^* \mid w \text{ endet auf } \word{10} \}  = \{\word{10}, \word{010}, \word{110}, \word{0010}, ...\}$
		\end{Beispiel}
\end{frame}







\begin{frame}{Produkt}
		\begin{Definition}
			Seien $L_1$ und $L_2$ zwei formale Sprachen. Dann bezeichnet
				$$L_1 \cdot L_2 := \{w_1 w_2 \mid w_1 \in L_1 \text{ und } w_2 \in L_2 \}$$
				das \textbf{Produkt} der Sprachen $L_1$ und $L_2$.
		\end{Definition}
		\pause
		In $L_1 \cdot L_2$ sind also alle Wörter enthalten, deren erster Teil aus $L_1$ und deren zweiter Teil aus $L_2$ ist.
	
\end{frame}

\begin{frame}{Produkt}
	\begin{Beispiele}
		$\{\word a, \word b\} \cdot \{\word c, \word d\} = \{\word{ac}, \word{ad}, \word{bc}, \word{bd}\}$\\[0.3em]
		\pause
		$ L = \{w_1w_2 \mid w_1\in \{\literal{b}\}^* ,  w_2\in \{\literal{a}\}^* \} 
		= \{\literal{b}\}^* \cdot \{\literal{a}\}^* $\\[1em]
		\pause
		Für alle formalen Sprachen $L$ gilt:\\
		$$ L \cdot \{\varepsilon\} = L  \qquad L \cdot \emptyset = \emptyset$$
	\end{Beispiele}
\end{frame}



\begin{frame}{Potenzen}
	\begin{Definition}
	Potenzen formaler Sprachen definieren wir induktiv so: 
	$$L^0 := \{\varepsilon \}$$
	$$L^{i+1} := L^i \cdot L$$
	\end{Definition} \pause
	$L^i$ enthält also alle Kombinationen von $i$ (nicht unbedingt verschiedenen) Wörtern aus $L$.
\end{frame}

\begin{frame}{Konkatenationsabschluss}
	\begin{Definition}
		Der \textbf{Konkatenationsabschluss} einer Sprache $L$ ist \quad $L^\ast = \bigcup \limits_{i=0}^\infty L^i$. \\
		$$ \set{}^* = \set{ \eps } \qquad \set{\word{abc}, \word{xy}}^* = \set{\eps, \word{abc}, \word{xy}, \word{abcabc}, \word{abcabcxy}, \word{xyabc}, ...} $$ 
		\medskip
		\pause
		Der \textbf{$\eps$-freie Konkatenationsabschluss} ist \quad $L^+ = \bigcup \limits_{i=1}^\infty L^i$. \\
		$$ \set{}^+ = \set{} \qquad \set{\word{abc}, \word{xy}}^+ = \set{\word{abc}, \word{xy}, \word{abcabc}, \word{abcabcxy}, \word{xyabc}, ...} $$
	\end{Definition} \pause
	\textbf{Achtung}: Der $\eps$-freie Konkatenationsabschluss muss nicht „$\eps$-frei“ sein! \\ \pause
	$$ \set{\eps}^+ = \underbrace{\set{\eps}}_{\mathclap{\text{enthält ein $\eps$!}}} $$
	
\end{frame}

\begin{frame}{Beispiele aus dem Leben}
	\begin{itemize}
		\item Sprache der korrekten IP4-Adressen: \pause
		\begin{align*}
		\set{\word{192.168.178.1}, \quad 
			\word{76.147.112.6}, ... }
		\end{align*} 
		\pause Enthält aber nicht: $\word{000.999.123.666}$ \pause
		\item Formale Sprache der Schlüsselwörter in Java: $$L = \{ \word{class}, \word{int} , \word{if}, \dots \}$$ \pause
		\item Formale Sprache der ganzen Zahlen: Mit $A = \{\word 0, ..., \word 9\}$ bauen wir
		\pause $$A \cdot A^* = \{\word{12849}, \word{1001}, \word{42}, ...\}$$ 
		\pause Und minus? Also besser $ \{ \mword-, \varepsilon \} \cdot A \cdot A^*$
	\end{itemize}
\end{frame}

\begin{frame}{Formale Sprachen}
	
	\begin{Beispiel}
		Sei $L$ formale Sprache aller Wörter über $A=\{\literal{a},\literal{b}\}$, in denen nirgends das Teilwort $\literal{ab}$ vorkommt.
		\begin{itemize}
			\pause
			\implitem $L=\{\literal{a},\literal{b}\}^*
			\setminus \set{w_1 \cdot \literal{ab} \cdot w_2 \mid w_1,w_2\in
			\{\literal{a},\literal{b}\}^* }$
			
			\pause
			\item Erst ein beliebiges Wort (evtl. $\varepsilon$) nur aus $\literal{b}$,\\
			danach ein beliebiges Wort (evtl. $\varepsilon$) nur aus $\literal{a}$.
			
			\pause
			\implitem $L=\set{w_1w_2 \mid w_1\in 
			\{\literal{b}\}^*  \text{ und }  w_2\in \{\literal{a}\}^* }$
		\end{itemize}
	\end{Beispiel}
	
	\begin{block}{Bemerkungen}
		\begin{itemize}
			\pause
			\item Die Beschreibung einer formalen Sprache ist \textbf{nicht} eindeutig (siehe oben).
			\pause
			\item Immer auf den Unterschied achten: $\{\literal{abc} \} \neq \literal{abc} $
		\end{itemize}
	\end{block}
\end{frame}


\begin{frame}{Aufgabe}
	Es sei $A = \{\word a, \word b\}$. Beschreibt die folgenden formalen Sprachen mit den Symbolen $\word\{, \word\}, \word a, \word b, \textcolor{blue}{\eps}, 
	\mword{\cup}, 
	\mword{{}^*}, \mword{{}^+},
	\word, , \mword{\cdot} , \word( \text{ und } 
	\word)
	$:
	\begin{itemize}
		\item die Menge aller Wörter über $A$, die das Teilwort \word{ab} enthalten. \\
		\uncover<2-|handout:2>{$\{\word a, \word b\}^\ast \cdot \{\word{ab}\} \cdot \{\word a, \word b\}^\ast$}
		\item die Menge aller Wörter über $A$, deren vorletztes Zeichen ein \word b ist. \\
		\uncover<3-|handout:2>{$\{\word a, \word b\}^\ast \cdot \{\word b\} \cdot \{\word a, \word b\}$}
		% $\{\word a, \word b\}^\ast \cdot \{\word b\} \cdot \{\word a, \word b\}$
		\item die Menge aller Wörter über $A$, in denen nirgends zwei \word bs unmittelbar hintereinander vorkommen. \\
		\uncover<4-|handout:2>{$\{\word a, \word{ba}\}^\ast \cdot \{ \word b, \eps \}$}
	\end{itemize}
\end{frame}




%\section{Binäre Operatoren}
\begin{frame}
	\begin{block}{Binäre Operatoren}
		Eine binäre Operation auf einer Menge $M$ ist eine Abbildung:
		\[ \diamond: M \times M \to M \]
		Üblicherweise benutzt man (wie bei manchen Relationen) die Infixschreibweise.
	\end{block} \pause
	\begin{block}{Eigenschaften}
		Die Operation $\diamond$ heißt kommutativ, wenn gilt:
		\[ \forall x \in M \ \forall y \in M: x \diamond y = y \diamond x \]
		
		Die Operation $\diamond$ heißt assoziativ, wenn gilt:
		\[ \forall x \in M \ \forall y \in M \ \forall z \in M : (x \diamond y) \diamond z = x \diamond (y \diamond z) \]
	\end{block} \pause


\end{frame}

\begin{frame}	
	\begin{block}{Was ihr nun wissen solltet}
		\begin{itemize}
			\item Aussagenlogik: Syntax und Semantik
			\item Was eine formale Sprache ist und warum das Konzept wichtig ist
			\item Wie man einfache formale Sprachen formal angeben kann
			\item Einfache Operationen auf formalen Sprachen
		\end{itemize}
	\end{block}
	
	\begin{block}{Was nächstes Mal kommt}
		\begin{itemize}
			\item So viele Sprachen: Mehr zu formalen Sprachen
			\item Aus 2 mach 10: Übersetzungen
		\end{itemize}
	\end{block}
\end{frame}

%% Letzte Seite
\lastframe{0.45}{0}{xkcd/logic_drivers.png}{https://xkcd.com/589/}
\slideThanks

\end{document}
